%\documentclass[a4paper, oneside, openany]{book}
\documentclass{article}

%**************************************************************
% Importazione package
%************************************************************** 

% permette di modificare i margini
%\usepackage[top=3.1cm, bottom=3.1cm, left=2.2cm, right=2.2cm]{geometry}
\usepackage[a4paper]{../../modello/geometry}

% specifica con quale codifica bisogna leggere i file
\usepackage[utf8]{inputenc}

% necessario per risolvere il problema del carattere invisibile per l'a capo
\DeclareUnicodeCharacter{00A0}{ }

% per scrivere in italiano e in inglese;
% l'ultima lingua (l'italiano) risulta predefinita
\usepackage[english, italian]{babel}

% imposta lo stile italiano per i paragrafi
\usepackage{parskip}

% fornice elenchi numerati in ordine inverso
\usepackage{etaremune}

% comandi per l'appendice
\usepackage{appendix}
\renewcommand\appendixtocname{Appendici}

% numera anche i paragrafi
\setcounter{secnumdepth}{4}

% elenca anche i paragrafi nell'indice
\setcounter{tocdepth}{4}

% permetti di definire dei colori
\usepackage[usenames,dvipsnames]{color}

% permette di usare il comando "paragraph" come subsubsubsection!
\usepackage{titlesec}

% permette di inserire le immagini/tabelle esattamente dove viene usato il
% comando \begin{figure}[H] ... \end{figure}
% evitando che venga spostato in automatico
\usepackage{float}

% permette l'inserimento di url e di altri tipi di collegamento
\usepackage[colorlinks=true]{hyperref}

\hypersetup{
    colorlinks=true, % false: boxed links; true: colored links
    citecolor=black,
    filecolor=black,
    linkcolor=black, % color of internal links
    urlcolor=Maroon  % color of external links
}

% permette al comando \url{...} di andare a capo a metà di un link
\usepackage{breakurl}

% immagini
\usepackage{graphicx}

% permette di riferirsi all'ultima pagina con \pageref{LastPage}
\usepackage{lastpage}

% tabelle su più pagine
\usepackage{longtable}

% per avere dei comandi in più da poter usare sulle tabelle
\usepackage{booktabs}

% tabelle con il campo X per riempire lo spazio rimanente sulla riga
\usepackage{tabularx}

% multirow per tabelle
\usepackage{multirow}

% Permette tabelle su più pagine
\usepackage{longtable}

% colore di sfondo per le celle
\usepackage[table]{xcolor}

% permette di fare longtable larghe tutta la pagina (parametro x)
% su Ubuntu non si può installare il pacchetto, deve essere in modello/
\usepackage{../../modello/tabu}

% imposta lo spazio tra le righe di una tabella
\setlength{\tabulinesep}{6pt}

% definisci un nuovo tipo di colonna P che permette di andare a capo con \newline
% e giustificata a sinistra
\usepackage{array}
\usepackage{ragged2e}
\newcolumntype{P}[1]{>{\RaggedRight \hspace{0pt}}p{#1}}

% personalizza l'intestazione e piè di pagina
\usepackage{fancyhdr}

% permette di inserire caratteri speciali
\usepackage{textcomp}

% permette di aggiustare i margini e centrare tabelle e figure
\usepackage{changepage}

%Permette di includere i grafici a barre
%IMPORTANTE: deve essere caricato prima di /pgfgantt altrimenti causa conflitto
\usepackage{pgfplots}

% permette di includere i diagrammi Gantt
% su Ubuntu non si può installare il pacchetto, deve essere in modello/
\usepackage{../../modello/pgfgantt}

% permette di includere i grafici a torta
\usepackage{../../modello/pgf-pie}

% necessario per pgf-pie
\usepackage{tikz}

% permette i path delle immagini con gli spazi
\usepackage{grffile}

% ruota le immagini
\usepackage{rotating}

% permetti di calcolare le larghezze facendo calcoli
\usepackage{calc}

\fancypagestyle{plain}{
	% cancella tutti i campi di intestazione e piè di pagina
	\fancyhf{}

	\lhead{
		\includegraphics[height=1.5cm, width=1.5cm, keepaspectratio=true]{../../modello/or-bit_bkg.png}
		\parbox[b]{10cm}{
			\emph{\GroupName{}} \vspace{0pt} \\
			\emph{Progetto \ProjectName{}} \vspace{7pt}
		}
	}
	\chead{}
	%\rhead{
	%	\slshape \leftmark
	%}

	\lfoot{
		\DocTitle{} \\
		% differenzia a seconda che \DocVersion{} stampi testo o no
		\setbox0=\hbox{\DocVersion{}\unskip}\ifdim\wd0=0pt
			% nulla
		\else
			v \DocVersion{}
		\fi
	}
	\rfoot{Pagina \thepage{} di \pageref{LastPage}}

	% Visualizza una linea orizzontale in cima e in fondo alla pagina
	\renewcommand{\headrulewidth}{0.3pt}
	\renewcommand{\footrulewidth}{0.3pt}
}
\setlength{\headheight}{30pt}
\pagestyle{plain}

% allarga l'header a tutta la pagina
%\fancyhfoffset[L]{\oddsidemargin + \hoffset + 1in}
%\fancyhfoffset[R]{\evensidemargin + \marginparwidth - \marginparsep}

% Per inserire del codice sorgente formattato
\usepackage{listings}

\lstset{
  extendedchars=true,          % lets you use non-ASCII characters
  inputencoding=utf8,   % converte i caratteri utf8 in latin1, richiede \usepackage{listingsutf8} anzichè listings
  basicstyle=\ttfamily,        % the size of the fonts that are used for the code
  breakatwhitespace=false,     % sets if automatic breaks should only happen at whitespace
  breaklines=true,             % sets automatic line breaking
  captionpos=t,                % sets the caption-position to top
  commentstyle=\color{mygreen},   % comment style
  frame=none,               % adds a frame around the code
  keepspaces=true,            % keeps spaces in text, useful for keeping indentation of code (possibly needs columns=flexible)
  keywordstyle=\bfseries,     % keyword style
  numbers=none,               % where to put the line-numbers; possible values are (none, left, right)
  numbersep=5pt,              % how far the line-numbers are from the code
  numberstyle=\color{mygray}, % the style that is used for the line-numbers
  rulecolor=\color{black},    % if not set, the frame-color may be changed on line-breaks within not-black text (e.g. comments (green here))
  showspaces=false,           % show spaces everywhere adding particular underscores; it overrides 'showstringspaces'
  showstringspaces=false,     % underline spaces within strings only
  showtabs=false,             % show tabs within strings adding particular underscores
  stepnumber=5,               % the step between two line-numbers. If it's 1, each line will be numbered
  stringstyle=\color{red},    % string literal style
  tabsize=4,                  % sets default tabsize
  firstnumber=1      % visualizza i numeri dalla prima linea
}

% Permetti di utilizzare il grassetto per i caratteri Typewriter (per es. il font di \code{...} e \file{...})
\usepackage[T1]{fontenc}
\usepackage{lmodern}

%%%%%%%%%%%%%%
%  COSTANTI  %
%%%%%%%%%%%%%%

% In questa prima parte vanno definite le 'costanti' utilizzate da due o più documenti.

% Meglio non mettere gli \emph dentro le costanti, in certi casi creano problemi
\newcommand{\GroupName}{Or-bit}
\newcommand{\GroupEmail}{orbit.swe@gmail.com}
\newcommand{\ProjectName}{Monolith}
\newcommand{\ProjectVersion}{v1.0.0}

\newcommand{\Proponente}{Red Babel}
\newcommand{\Committente}{Prof. Tullio Vardanega \\ Prof. Riccardo Cardin}
\newcommand{\ResponsabileInCarica}{}

% La versione dei documenti deve essere definita qui in global, perchè serve anche agli altri documenti
\newcommand{\VersioneG}{1.0.0}
\newcommand{\VersionePQ}{1.0.0}
\newcommand{\VersioneNP}{1.0.0}
\newcommand{\VersionePP}{1.0.0}
\newcommand{\VersioneAR}{1.0.0}
\newcommand{\VersioneSF}{1.0.0}
\newcommand{\VersioneST}{0.0.0}
\newcommand{\VersioneMA}{0.0.0}
\newcommand{\VersioneMS}{0.0.0}
\newcommand{\VersioneMU}{0.0.0}
\newcommand{\VersioneDP}{0.0.0}
% Il verbale non ha versionamento.
% Lasciare vuoto, non mettere trattini o puntini
% Non sono permessi numeri nel nome di un comando :(
\newcommand{\VersioneVprimo}{} 
\newcommand{\VersioneVsecondo}{}

%ATTENZIONE: vengono rimossi dai comandi sottostanti i numeri di versione dei documenti non ancora esistenti
% Quando serve riferirsi a ``Nome del Documento + ultima versione x.y.z'' usiamo queste costanti:
\newcommand{\Glossario}{\emph{Glossario v\VersioneG{}}}
\newcommand{\PianoDiQualifica}{\emph{Piano di Qualifica v\VersionePQ{}}}
\newcommand{\NormeDiProgetto}{\emph{Norme di Progetto v\VersioneNP{}}}
\newcommand{\PianoDiProgetto}{\emph{Piano di Progetto v\VersionePP{}}}
\newcommand{\StudioDiFattibilita}{\emph{Studio di Fattibilità v\VersioneSF{}}}
\newcommand{\AnalisiDeiRequisiti}{\emph{Analisi dei Requisiti v\VersioneAR{}}}
\newcommand{\SpecificaTecnica}{\emph{Specifica Tecnica}}
\newcommand{\ManualeAdmin}{\emph{Manuale Admin}}
\newcommand{\ManualeSviluppatore}{\emph{Manuale Sviluppatore}}
\newcommand{\ManualeUtente}{\emph{Manuale Utente}}
\newcommand{\DefinizioneDiProdotto}{\emph{Definizione di Prodotto}}

% Per riferirsi alle revisioni usare i comandi seguenti:
\newcommand{\RR}{Revisione dei Requisiti}
\newcommand{\RP}{Revisione di Progettazione}
\newcommand{\RQ}{Revisione di Qualifica}
\newcommand{\RA}{Revisione di Accettazione}

% Comandi utili per riferirsi alle varie fasi
\newcommand{\AR}{Analisi dei Requisiti}
\newcommand{\AD}{Analisi di Dettaglio}
\newcommand{\PA}{Progettazione Architetturale}
\newcommand{\PD}{Progettazione di Dettaglio}
\newcommand{\Cod}{Codifica}
\newcommand{\VV}{Verifica e Validazione}

% Comandi utili per riferirsi ai ruoli di progetto
\newcommand{\Responsabile}{\textit{Responsabile di Progetto}}
\newcommand{\Amministratore}{\textit{Amministratore}}
\newcommand{\Analista}{\textit{Analista}}
\newcommand{\Progettista}{\textit{Progettista}}
\newcommand{\Programmatore}{\textit{Programmatore}}
\newcommand{\Verificatore}{\textit{Verificatore}}

\newcommand{\Responsabili}{\textit{Responsabili di Progetto}}
\newcommand{\Amministratori}{\textit{Amministratori}}
\newcommand{\Analisti}{\textit{Analisti}}
\newcommand{\Progettisti}{\textit{Progettisti}}
\newcommand{\Programmatori}{\textit{Programmatori}}
\newcommand{\Verificatori}{\textit{Verificatori}}


\newcommand{\ScopoDelProdotto}{
Lo scopo di questo progetto è sviluppare un framework volto alla costruzione di bubble interattive, per aggiungere funzionalità al programma di comunicazione Rocket.Chat.
In secondo luogo, utilizzare il framework prodotto per fornire esempi significativi delle suddette bubble.
}

\newcommand{\GlossarioIntroduzione}{
Insieme al progetto il gruppo fornirà anche un glossario nel quale saranno raccolti i termini tecnici, acronimi e le parole che necessitano di chiarimenti, che verranno utilizzati nella documentazione.\\
I vocaboli presenti nel documento \Glossario{} saranno marcati con una \textit{\small{G}} al pedice.
}

% Altri comandi utili
\newcommand{\email}{e-mail}
\newcommand{\sezione}{sezione §}

%%%%%%%%%%%%%%
%  FUNZIONI  %
%%%%%%%%%%%%%%

% In questa seconda parte vanno definite le 'funzioni' utilizzate da due o più documenti.

% Comando per usare comodamente le virgolette
\newcommand{\virgolette}[1]{``#1''}

% Serve a dare la giusta formattazione alle parole presenti nel glossario
% il nome del comando \glossary è già usato da LaTeX
\newcommand{\glossario}[1]{\textit{#1\ped{\ped{G}}}}

\newcommand*\flextt{%
  \fontdimen2\font=0.4em% interword space
  \fontdimen3\font=0.6em% interword stretch
  \fontdimen4\font=0.1em% interword shrink
  \fontdimen7\font=0.1em% extra space
  \hyphenchar\font=`\:
}

% Serve a dare la giusta formattazione al codice inline
\newcommand{\code}[1]{\flextt{\texttt{#1}}}

% Serve a dare la giusta formattazione a tutte le path presenti nei documenti
\newcommand{\file}[1]{\flextt{\texttt{#1}}}

% Permette di andare a capo all'interno di una cella in una tabella
\newcommand{\multiLineCell}[2][c]{\begin{tabular}[#1]{@{}l@{}}#2\end{tabular}}

% Genera automaticamente la pagina di copertina
\newcommand{\makeFrontPage}{
  % Declare new goemetry for the title page only.
  \newgeometry{top=3.5cm}
  
  \begin{titlepage}
  \begin{center}

  \begin{center}
  \includegraphics[width=8cm]{../../modello/or-bit_bkg.png}
  \end{center}
  
  \vspace{1cm}

  \begin{Huge}
  \textbf{\DocTitle{}}
  \end{Huge}
  
  \textbf{\emph{Gruppo \GroupName{} \, \texttwelveudash{} \, Progetto \ProjectName{}}}
  
  \vspace{11pt}

  \bgroup
  \def\arraystretch{1.3}
  \begin{tabular}{ r|l }
    \multicolumn{2}{c}{\textbf{Informazioni sul documento}} \\
    \hline
		% differenzia a seconda che \DocVersion{} stampi testo o no
		\setbox0=\hbox{\DocVersion{}\unskip}\ifdim\wd0=0pt
			% nulla (non ho trovato come togliere l'a capo)
			\\
		\else
			\textbf{Versione} & \DocVersion{} \\
		\fi
    \textbf{Redazione} & \multiLineCell[t]{\DocRedazione{}} \\
    \textbf{Verifica} & \multiLineCell[t]{\DocVerifica{}} \\
    \textbf{Approvazione} & \multiLineCell[t]{\DocApprovazione{}} \\
    \textbf{Uso} & \DocUso{} \\
    \textbf{Distribuzione} & \multiLineCell[t]{\DocDistribuzione{}} \\
  \end{tabular}
  \egroup

  \vspace{22pt}

  \textbf{Descrizione} \\
  \DocDescription{}

  \end{center}
  \end{titlepage}
  
  % Ends the declared geometry for the titlepage
  \restoregeometry
}
%%%%%%%%%%%%%%
%  COSTANTI  %
%%%%%%%%%%%%%%
\usepackage{amsmath}

% In questa prima parte vanno definite le 'costanti' utilizzate soltanto da questo documento.
% Devono iniziare con una lettera maiuscola per distinguersi dalle funzioni.

\newcommand{\DocTitle}{Piano di Qualifica}
\newcommand{\DocVersion}{\VersionePQ{}}

\newcommand{\DocRedazione}{}
\newcommand{\DocVerifica}{}
\newcommand{\DocApprovazione}{}

\newcommand{\DocUso}{Esterno}
\newcommand{\DocDistribuzione}{
	\Committente{} \\
	Gruppo \GroupName{}
}

% La descrizione del documento
\newcommand{\DocDescription}{
 Documento riguardante le strategie di verifica e gli obiettivi qualitativi per il progetto \ProjectName{} adottate dal gruppo \GroupName.
 }

%%%%%%%%%%%%%%
%  FUNZIONI  %
%%%%%%%%%%%%%%

% In questa seconda parte vanno definite le 'funzioni' utilizzate soltanto da questo documento.


\begin{document}

\makeFrontPage

\section*{Diario delle Modifiche}


\bgroup
\begin{longtable}{|c|c|c|c|P{4.2cm}|}
	\hline \multicolumn{1}{|c|}{\textbf{Versione}} & \multicolumn{1}{c|}{\textbf{Data}} & \multicolumn{1}{c|}{\textbf{Persona}} & \multicolumn{1}{c|}{\textbf{Ruolo}} & \multicolumn{1}{c|}{\textbf{Descrizione}} \\ \hline 
	\endfirsthead
	
	\hline \multicolumn{1}{|c|}{\textbf{Versione}} & \multicolumn{1}{c|}{\textbf{Data}} & \multicolumn{1}{c|}{\textbf{Persona}} & \multicolumn{1}{c|}{\textbf{Ruolo}} & \multicolumn{1}{c|}{\textbf{Descrizione}} \\ \hline 
	\endhead
	
	\hline \multicolumn{5}{|r|}{{Continua nella prossima pagina}} \\ \hline
	\endfoot
	
	\hline \hline
	\endlastfoot
	
	% IN ORDINE DALLA MODIFICA PIÙ RECENTE ALLA PIÙ VECCHIA
	 \hline 1.0.0 & 17-12-2016 & Mattia Zecchinato & Responsabile & Approvazione del documento \\
	 
	 \hline 0.2.0 & 16-12-2016 & Marco Ponchia & Verificatore & Seconda verifica documento \\
	 
	 \hline 0.1.0 & 15-12-2016 & Beatrice Guerra & Verificatore & Prima verifica del documento \\
	 
	 \hline 0.0.8 & 14-12-2016 & Tommaso Rosso & Amministratore & Completamento sezione \virgolette{Procedure a supporto dei processi} \\
	 
	 \hline 0.0.7 & 14-12-2016 & Alice V. Sasso & Amministratore & Completamento sezione \virgolette{Ambiente di lavoro} \\
	 
	 \hline 0.0.6 & 13-12-2016 & Alice V. Sasso & Amministratore & Inizio stesura sezione \virgolette{Ambiente di lavoro} \\
	 
	 \hline 0.0.5 & 12-12-2016 & Lorenzo Ferrarin & Amministratore & Riorganizzazione sezione \virgolette{Riunioni} \\
	 
	 \hline 0.0.4 & 12-12-2016 & Nicola Dal Maso & Amministratore & Stesura sezione \virgolette{Procedure a supporto dei processi} \\
	 
	 \hline 0.0.4 & 11-12-2016 & Nicola Dal Maso & Amministratore & Stesura sezione \virgolette{Ruoli di progetto} \\
	 
	 \hline 0.0.3 & 10-12-2016 & Alice V. Sasso & Amministratore & Stesura sezione \virgolette{Documentazione} \\
	  
	 \hline 0.0.2 & 10-12-2016 & Lorenzo Ferrarin & Amministratore & Stesura \virgolette{Collaborazione} \\
	
	 \hline 0.0.1 & 09-12-2016 & Lorenzo Ferrarin & Amministratore & Creazione struttura e stesura \virgolette{Introduzione} \\
	
	 \hline
\end{longtable}
\egroup


\clearpage
\tableofcontents

\letteraGlossario{A}
\definizione{Actor Model}
Modello matematico di esecuzione concorrente di un programma nel quale le primitive di elaborazione concorrente sono individuate negli attori.

\definizione{Akka}
Libreria \glossario{open source} che semplifica la costruzione di applicazioni concorrenti e distribuite per \glossario{JVM}.

\definizione{Algoritmo Predittivo}
Algoritmo che cerca di \virgolette{fare previsioni} sugli eventi futuri in base a dati raccolti in eventi passati.

\definizione{Amazon}
Societ\`a di commercio elettronico statunitense con sede a Seattle.

\definizione{Amazon web service}
 Raccolta di servizi web remoti offerti da \glossario{Amazon} operanti da 11 regioni geografiche distinte.

\definizione{Apache}
Organizzazione no-profit americana a sostegno di progetti software.

\definizione{API}
Acronimo di application programming interface, rappresentano le librerie software disponibili per un certo linguaggio di programmazione.
\clearpage

\letteraGlossario{B}
\definizione{Beacons}
Piccoli dispositivi che, attraverso la tecnologia \glossario{BLE}, sono in grado di trasmettere informazioni a smartphone e tablet.

\definizione{BLE}
Acronimo di Bluetooth Low Energy \`e una tecnologia ti tipo \glossario{Bluetooth} a basso consumo energetico.

\definizione{Bluetooth}
Standard di trasmissione per reti personali senza fili, sicuro ed economico che utilizza una frequenza radio a corto raggio.

\definizione{Bug}
Identifica un errore nella scrittura di un programma software.

\definizione{Bug Tracking}
Applicativo software utile al team di sviluppo di un progetto per tenere traccia delle segnalazioni di \glossario{bug} trovati nel proprio prodotto.
 
\definizione{Bytecode}
Linguaggio intermedio tra linguaggio di programmazione e linguaggio macchina che riduce l'indipendenza dall'hardware.
\clearpage

\letteraGlossario{C}
\definizione{Case}
Si accompagna agli aggettivi \textit{lower} o \textit{upper}, sta ad indicare rispettivamente se una lettera è minuscola o maiuscola.

\definizione{Caso d'uso}
Funzionalità di un prodotto software, o tecnica usata nei processi di ingegneria del software per effettuare in maniera esaustiva e non ambigua la raccolta dei requisiti, al fine di produrre software di qualità.

\definizione{Commit}
Parlando di \glossario{controllo di versione}, un commit si effettua quando si copiano le modifiche fatte su file locali nella cartella del \glossario{repository}.

\definizione{Controllo di Versione}
Gestione di versioni multiple di un insieme di informazioni, siano questi documenti testuali o parti di un programma software.

\definizione{CSS}
Acronimo di Cascading Style Sheets, ovvero fogli di stile, \`e un linguaggio usato nella formattazione di pagine web.

\definizione{CSS3}
Specifiche riguardati i fogli di stile \glossario{CSS}, costituite da sezioni separate dette moduli, con differenti stati di avanzamento e stabilit\`a. Alcuni di questi moduli sono stati pubblicati formalmente come \glossario{W3C Recommendation}, nel novembre 2014.
\clearpage

\letteraGlossario{D}
\definizione{Design Pattern}
Concetto che può essere definito come una soluzione progettuale generale ad un problema ricorrente. Si tratta di una descrizione o modello logico da applicare per la risoluzione di un problema che può presentarsi in diverse situazioni durante le fasi di progettazione e sviluppo del software, ancor prima della definizione dell'algoritmo risolutivo della parte computazionale.

\definizione{Debugging}
Attività che consiste nell'individuazione da parte del programmatore della porzione di software affetta da errore (\glossario{bug}) rilevata nei software a seguito dell'utilizzo del programma.
\clearpage

\letteraGlossario{F}
\definizione{Framework}
Architettura logica di supporto (spesso un'implementazione logica di un particolare \glossario{design pattern}) su cui un software può essere progettato e realizzato, spesso facilitandone lo sviluppo da parte del team di sviluppatori.
\clearpage

\letteraGlossario{G}
\definizione{Gantt}
Ideatore del \textit{diagramma di Gantt}, strumento usato nelle attività di \glossario{project management} per tenere sotto controllo tutte le attività correlate al progetto in una determinata fascia temporale.

\definizione{Git}
Sistema software di \glossario{controllo di versione} distribuito.

\definizione{GitHub}
GitHub è un servizio web di hosting per lo sviluppo di progetti software, che usa il sistema di controllo di versione \glossario{Git}.
\clearpage

\letteraGlossario{H}
\definizione{HTML}
Linguaggio di markup per la strutturazione di pagine web.

\definizione{HTML5}
Linguaggio \glossario{HTML} pubblicato come \glossario{W3C Recommendation} da ottobre 2014.

\definizione{HTTP}
Acronimo di HyperText Transfer Protocol \`e un protocollo utilizzato come sistema per la trasmissione di informazioni sul web. Le specifiche del protocollo sono gestite dal World Wide Web Consortium (W3C).

\definizione{HTTPS}
Acronimo di  HyperText Transfer Protocol over Secure Socket Layer \`e l'applicazione della crittografia asimmetrica al protocollo \glossario{HTTP}, per garantire trasferimenti di dati nel web evitando attacchi di tipo \glossario{man in the middle}.
\clearpage

\letteraGlossario{I}
\definizione{IDE}
Acronimo per Integrated Development Environment, ovvero un ambiente di sviluppo integrato per la realizzazione di programmi per sistemi informatici.

\definizione{Instant Messaging}
Categoria di sistemi di comunicazione in tempo reale in rete, tipicamente Internet o una rete locale, che permette ai suoi utilizzatori lo scambio di brevi messaggi.

\definizione{Internet of Things}
Neologismo riferito all'estensione di Internet al mondo degli oggetti e dei luoghi concreti.

\definizione{IoT}
Acronimo di \glossario{Internet of Things}.

\definizione{ISO}
Abbreviazione per International Organization for Standardization, è la più importante organizzazione a livello mondiale per la definizione di norme tecniche.
\clearpage

\letteraGlossario{J}
\definizione{Java}
linguaggio di programmazione orientato agli oggetti, progettato per essere indipendente dalla piattaforma di esecuzione, utilizzando l'implementazione di un processore virtuale detto \glossario{JVM}.

\definizione{Javascript}
Linguaggio di scripting orientato agli oggetti e agli eventi, utilizzato nella programmazione Web lato client per la creazione di effetti dinamici n siti e applicazioni web.

\definizione{JSON}
Acronimo di JavaScript Object Notation, \`e un formato usato nell''interscambio di dati fra applicazioni client-server.

\definizione{JVM}
Acronimo di Java Virtual Machine \`e il componente della piattaforma Java che esegue i programmi tradotti in \glossario{bytecode} dopo una prima compilazione.
\clearpage

\letteraGlossario{M}
\definizione{Machine learning}
Rappresenta un'area fondamentale nell'intelligenza artificiale e realizzando algoritmi e sistemi basti su osservazioni di dati per la sintesi di nuova conoscenza.

\definizione{Man in the middle}
Tipo di attacco crittografico nel quale l'attaccante \`e in grado di leggere, inserire o modificare a piacere, messaggi tra due parti comunicanti tra di loro.

\definizione{Milestone}
Importante traguardo intermedio nello svolgimento del progetto. Molto spesso è rappresentata da eventi, cioè da attività con durata zero o di un giorno, e viene evidenziata in maniera diversa dalle altre attività nell'ambito dei documenti di progetto. Può essere intesa anche come una particolare configurazione di item relativi al progetto.

\definizione{MongoDB}
Database non relazionale classificato come \glossario{NoSQL}, orientato ai documenti si tratta di software libero e \glossario{open source}.

\definizione{MQTT}
Acronimo di MQ Telemetry Transport (MQTT) \`e un protocollo di messaggistica posizionato in cima a TCP/IP, per situazioni a basso impatto e con banda limitata, che interagisce con un message broker responsabile della distribuzione dei messaggi ai client destinatari.
\clearpage

\letteraGlossario{N}
\definizione{Node.js}
\glossario{Framework} per lo sviluppo di applicazioni server-side di Javascript

\definizione{NoSQL}
Acronimo di Not Only SQL rappresentano tutti cui software che non utilizzano il modello relazionale, definendo le basi di date costruite in questo modo come memorizzazioni strutturate.
\clearpage

\letteraGlossario{O}
\definizione{Open source}
Accostato ad un software sta ad indicare che il codice sorgente dello stesso \`e pubblico, favorendone lo studio, le modifiche ed estensioni da parte di programmatori indipendenti.
\clearpage

\letteraGlossario{P}
\definizione{Package}
Collezione di classi e interfacce correlate.

\definizione{Play Framework}
\glossario{Framework} \glossario{open source} scritto in \glossario{Java} e \glossario{Scala}, ad elevata produttivit\`a che integra i componenti e le \glossario{API} necessarie per un moderno sviluppo delle applicazioni web.

\definizione{PNG}
Sigla per Portable Network Graphics, ovvero un formato file per immagini.

\definizione{Project Management}
Insieme di attività svolte tipicamene da una figura dedicata e specializzata detta project manager, volte all'analisi, alla progettazione, alla pianificazione e alla realizzazione degli obiettivi di un progetto, gestendolo in tutte le sue caratteristiche e fasi evolutive, nel rispetto di precisi vincoli (tempi, costi, risorse, scopi, qualità).
\clearpage

\letteraGlossario{R}
\definizione{Repository}
Luogo di memorizzazione dei file, spesso situato in un server remoto.

\definizione{Revert}
In ambito di \glossario{controllo di versione}, è l'abbandono di uno o più cambiamenti recenti in favore di un ritorno ad una precedente versione di un documento o di parti di software.
\clearpage

\letteraGlossario{S}
\definizione{Scala}
linguaggio di programmazione di tipo general-purpose multi-paradigma studiato per integrare le caratteristiche e funzionalit\`a dei linguaggi orientati agli oggetti e dei linguaggi funzionali.

\definizione{Skype}
Software proprietario freeware di \glossario{instant messaging} e VoIP. Con esso sono possibili le videochiamate e lo scambio di messaggi testuali o di file.

\definizione{SQL}
Acronimo di Structured Query Language \`e un linguaggio standardizzato per database che utilizzano il modello relazionale.

\definizione{SVG}
Acronimo di Scalable Vector Graphics, indica una tecnologia in grado di visualizzare oggetti di grafica vettoriale e, pertanto, di gestire immagini scalabili dimensionalmente.
\clearpage

\letteraGlossario{T}
\definizione{Telegram}
Servizio di \glossario{instant messaging} multipiattaforma, usato per inviare messaggi tramite Internet.

\definizione{Tomcat}
Server Web \glossario{open source} sviluppato dalla \glossario{Apache}. 

\definizione{Twitter Boostrap}
Raccolta di strumenti per la creazione di siti web, contenente modelli basati su \glossario{HTML}, \glossario{CSS} e \glossario{Javascript}.
\clearpage

\letteraGlossario{U}
\definizione{UML}
Acronimo per Unified Modeling Language, è un linguaggio visuale di modellazione e specifica basato sul paradigma object-oriented.

\definizione{UTF-8}
UTF-8, ovvero Unicode Transformation Format, 8 bit, è una codifica dei caratteri Unicode in sequenze di lunghezza variabile di byte.
\clearpage

\letteraGlossario{W}
\definizione{Way of Working}
La prassi, il modo di fare regolato da norme che segue un team o un'azienda nella produzione di prodotti software.

\definizione{W3C}
Sintesi di World Wide Web Consortium \`e la principale organizzazione per gli standard del World Wide Web.

\definizione{W3C Recommendation}
Standard formalmente dichiarati da parte del \glossario{W3C}.
\clearpage

\end{document}
